\documentclass[conference]{IEEEtran}
\usepackage[utf8]{inputenc}
\usepackage[T1]{fontenc}
\usepackage{cite}
\usepackage{amsmath,amssymb}
\usepackage{graphicx}
\usepackage{booktabs}
\usepackage{url}

\title{Verifier-Gated LLM NetOps: A Preregistered, Artifact-Backed Evaluation of Input-Sanitization and Deterministic Verification Against Adversarial Intents}

\author{
\IEEEauthorblockN{Autonomous AI Researcher}
\IEEEauthorblockA{CNSM 2026 Agentic AI Researcher Track}
}

\begin{document}

\maketitle

\begin{abstract}
We report a preregistered paired confirmatory experiment (study\_id cand-2026-03) that measured whether template-based input sanitization combined with deterministic verifier-gating (and at-most-one automated repair) changes the rate at which a deterministic validator accepts LLM-generated CLI configurations under adversarially mutated intents. The frozen run used the declared generator gpt-5-mini and deterministic\_netops\_validator\_v5 within the hosted\_netops\_gvr\_v1 adapter. Planned episodes = 300; completed episodes = 282; complete paired outcomes used in the primary analysis = 141. Observed per-pair acceptance: Baseline 135/141 (95.7447\%); Guarded 141/141 (100.0\%); paired absolute difference = 0.04255319148936165; exact two-sided McNemar p = 0.03125; 95\% bootstrap percentile CI (10,000 resamples, seed = 1729) = [0.014184397163120567, 0.07801418439716312]. We reconcile preregistration language vs the formal estimand, document per-condition pipeline semantics with explicit archived-artifact references, report the protocol deviation (unset runtime sampling temperature), and discuss operational interpretation and limitations constrained by synthetic scope and single-model/validator evaluation. All principal numeric claims are supported by archived execution and analysis artifacts referenced in Methods and Results.
\end{abstract}

\section{Introduction}
Generative LLMs are increasingly proposed to translate high-level operator intent into device CLI, promising automation gains for NetOps. However, operational risks remain when generated outputs violate sequencing constraints, CLI grammar, or safety invariants; adversarially crafted inputs can exacerbate these risks. Prior prototype studies show that coupling LLM generation with deterministic verification and targeted repair can reduce observable errors, but existing evaluations often rely on token- or reference-match metrics that do not directly measure deployability or acceptance by a deterministic validator [3,8,9].

This work asks a narrowly scoped, operational question amenable to preregistered inference: How effective are template-based input sanitization and deterministic verifier-gating (with at-most-one automated repair) at changing the rate at which a deterministic validator accepts LLM-generated CLI configurations when inputs are adversarially mutated? The study cand-2026-03 preregistered a single formal estimand (paired\_success\_rate\_difference\_guarded\_minus\_baseline) and a confirmatory paired test. Crucially, because the registered generation semantics were shared\_initial\_candidate, the paired comparison isolates downstream guarded-pipeline effects (sanitization, gating, permitted repair, and final validation) applied to a single shared LLM draft per intent rather than differences in initial generation.

Contributions. (1) A preregistered, artifact-backed paired experiment executed end-to-end and reconciled against the archived preregistration; (2) a precise, artifact-grounded description of per-condition pipeline semantics that demonstrates shared\_initial\_candidate reuse; (3) quantitative confirmatory results showing a small but statistically detectable absolute change in deterministic-validator acceptance under the guarded pipeline on a synthetic adversarial-intent bank; and (4) an explicit reconciliation of preregistration natural-language hypothesis wording with the formal estimand and a disclosure of a protocol deviation (unset sampling temperature) together with reproducibility guidance.

\section{Related Work}
Deterministic network-verification tools are a practical foundation for validator-gated NetOps pipelines because they check forwarding and configuration invariants at the configuration or header-space level and can run fast enough for interactive uses. Classic header-space analyses and their derivatives provide reachability and invariant-checking semantics over packet headers and device-wide rule sets, while streaming/online tools (e.g., VeriFlow/NetPlumber) focus on real-time policy checking as configurations change [W1882012874,W158224344,W2122695394]. These verifiers are valuable for gating because they give binary deployability decisions (accept/reject) against explicit formal properties; their effectiveness, however, depends on the granularity and correctness of the formal specification (what is asserted as a safety invariant) and on whether the verifier models dynamic protocol-level behavior or only static configuration-level constraints.

Empirical work on LLM-assisted configuration highlights complementary limitations and opportunities. Evaluations such as NetConfEval and other recent NetOps-focused studies report that LLMs frequently produce syntactic, ordering, or semantic-sequence errors when tasked with intent-to-CLI translation; generate--validate--repair architectures---where an LLM draft is checked by a deterministic tool and then repaired---have been demonstrated in prototype settings to reduce observable errors compared to unconstrained generation [W4399629579,W4388629193]. Survey and review articles in AIOps and telecommunications emphasize fragmented evaluation practices and call for operationally meaningful metrics (deployability, validator acceptance, semantic safety) rather than only token-level similarity or BLEU-style scores [W4411735779,W4402742293,W4383605161]. Those syntheses also note that many demonstrations remain prototype-scale and that improved benchmarks and measurement semantics are needed to compare mitigation strategies reliably.

On the security and orchestration side, prompt-injection and indirect-prompt attacks against LLM-integrated applications have been documented and motivate defenses such as input-sanitization, policy gating, and stricter tool-mediated checks in operational pipelines [W4388886073,W4407163436]. Concurrent work on tool-augmented LLM patterns (e.g., Toolformer-style approaches) shows that LLMs can be taught or orchestrated to call external checks and then use feedback to repair outputs, but these approaches depend on the fidelity of the external checks and the allowed repair scope; imperfect verifier expressiveness or permissive repair policies can convert otherwise-rejectable drafts into accepted ones without guaranteeing semantic-preservation [W4320165837,W4320165837].

Positioning. Building on these bodies of work, our preregistered experiment measures a validator-centered operational outcome (deterministic-validator acceptance) under a controlled, paired design that isolates post-generation pipeline interventions. The experiment and its artifacts explicitly tie LLM draft reuse, sanitizer application, gating decision, and any limited repair calls to archived traces so that the observed paired difference can be interpreted relative to verifier semantics and prior prototype findings rather than against token-similarity baselines.

\section{Methodology}
Overview and preregistration. The study was preregistered as cand-2026-03 before observing outcomes (preregistration\_sha256: ba6ed25e\allowbreak{}84c7240f\allowbreak{}be94f2d5\allowbreak{}0b3e778f\allowbreak{}bfffd810\allowbreak{}ff2db632\allowbreak{}df8946e8\allowbreak{}5f0f1373). The preregistration specifies a paired confirmatory design (shared\_initial\_candidate generation semantics), a single primary estimand (paired\_success\_rate\_difference\_guarded\_minus\_baseline), and concrete analysis procedures (exact two-sided McNemar for the paired comparison; bootstrap-percentile 95\% CIs with 10,000 resamples and seed = 1729). These analysis parameters and bootstrap metadata are archived (analysis/analysis\_log.jsonl).

Task bank, adversary model, and scope. Tasks were produced by deterministic\_netops\_task\_generator\_v5; each task encodes: (i) an initial device-state snapshot, (ii) an operator intent template (intent-to-configuration), and (iii) a required command sequence and per-field constraints (e.g., must-shutdown-before-MTU-change). The synthetic bank exercises admin, MTU, and VLAN updates and a small set of workflow policies (protected interfaces, group-activity minima, must-be-down constraints). The adversary model used deterministic mutation heuristics applied to these templates (string-level mutations, slot-overrides, token insertions) as archived in execution/task\_manifest.jsonl (see representative entries such as execution/task\_manifest.jsonl -> task-000001 payload). Scope was intentionally limited to these synthetic templates and the deterministic CLI grammar encoded in the validator; no production device logs or vendor-specific private data were used.

Orchestration, declared models, and protocol deviation. The orchestration adapter was hosted\_netops\_gvr\_v1 and the declared generator was gpt-5-mini (execution/model\_configuration.json). The deterministic validator was deterministic\_netops\_validator\_v5; validator outputs include explicit validation\_before and validation\_after objects with parsed command counts, touched-field counts, enumerated violation\_codes, and final\_state snapshots (examples: execution/scoring/task-000001-baseline.json and execution/scoring/task-000001-guarded.json). The preregistered pipeline constrained generation semantics to shared\_initial\_candidate (independent\_condition\_generation = false). The archived model\_configuration.json records declared\_model\_version = "gpt-5-mini" and temperature = null; we disclose this unset temperature as a protocol deviation because provider-side defaults may influence sampling. To preserve internal comparability despite the unset parameter the pipeline used a single shared generation per pair (evidence: analysis/deterministic\_reconciliation.json: stage\_counts.shared\_initial\_generation = 150) and provider-call accounting (hosted\_model\_call\_event\_count = 156, completed = 147, failed = 9) corroborates the actual call pattern.

Per-pair timeline and concrete artifact evidence (shared\_initial\_candidate semantics). The registered shared\_initial\_candidate semantics produced the following deterministic per-pair timeline; each step and its archived artifact(s) are listed to permit exact forensic inspection:
1) Task instantiation and adversarial mutation: recorded in execution/task\_manifest.jsonl (task payload contains the sanitized-orig intent and any slot changes). Representative entry: task-000001 payload (execution/model\_manifest and execution/task\_manifest.jsonl). 
2) Shared initial generation (single LLM call per pair): output saved to execution/responses/task-<id>-shared-initial.txt (e.g., execution/responses/task-000001-shared-initial.txt). Provider-stage accounting records shared\_initial\_generation = 150 (analysis/deterministic\_reconciliation.json).
3) Baseline branch scoring: deterministic\_netops\_validator\_v5 evaluates the shared\_initial\_candidate; baseline validator trace is in execution/scoring/task-<id>-baseline.json (validation\_before and validation\_after objects show parsed\_command\_count, required\_command\_count, valid flag, and violation codes when present).
4) Guarded branch pre-processing and gating: a template-based input-sanitizer is applied (sanitizer logic and results are recorded in the task payload fields). The sanitized candidate is presented to the verifier-gate; if the sanitized candidate fails gating the policy allows at-most-one automated repair call. Repair-stage provider-call traces for episodes invoking repair are present under execution/provider\_calls and referenced by execution/scoring artifacts (repair\_provider\_trace\_path fields). The pipeline-level repair budget (<=1 repair call per episode) is enforced by the orchestration adapter (preregistration and execution manifests). Evidence: stage\_counts.repair = 6 and repair traces linked in execution/scoring files and in analysis/deterministic\_reconciliation.json.
5) Final guarded validation: the validator re-evaluates the sanitized or repaired candidate, storing validation\_after and final\_configuration in execution/scoring/task-<id>-guarded.json. Paired artifacts for the same shared\_initial\_candidate include the shared\_initial\_candidate\_sha256 value present in both baseline and guarded scoring artifacts (e.g., execution/scoring/task-000001-baseline.json and execution/scoring/task-000001-guarded.json share the same shared\_initial\_candidate\_sha256).

Sanitizer and gating specifics (artifact-grounded). The input-sanitizer is template-driven: it enforces required slots, normalizes numeric fields (e.g., MTU integer normalization), and rejects or rewrites fields that violate per-task allowed\_fields/allowed\_touched\_fields constraints; sanitized outputs are recorded in the task payload (execution/task\_manifest.jsonl task\_payload.constraints and task\_payload.intent). The gating policy is deterministic: if deterministic\_netops\_validator\_v5 reports violations that fall into a policy-rejection set (e.g., protected-interface modification or missing required-step violations), the candidate is rejected; if violations are repairable under permitted-scope rules, the orchestrator may invoke a single repair generation attempt. These behaviors and the task-specific policies are visible in the task payload constraints and in per-episode validator traces (execution/scoring/* files).

Repair semantics and auditability. Repairs are restricted by preregistered policy: one additional LLM call per episode, limited to localized edits (slot corrections, sequence reordering) rather than free-form regeneration. Repair calls and resulting candidates are fully auditable: provider traces, shared\_initial\_candidate\_sha256, repair\_prompt\_sha256, and repair\_provider\_trace\_path appear in execution/scoring artifacts for repaired episodes (representative repaired episodes are traceable and correspond to the six repair-stage counts in analysis/deterministic\_reconciliation.json). The six discordant n10 cases observed in the paired contingency table correspond to episodes where a permitted repair converted an initially rejected draft into a validator-accepted final candidate; individual scoring artifacts show validation\_before (invalid) and validation\_after (valid) snapshots.

Outcome definitions, missingness handling, and exclusions. For each adversarial intent (paired unit) the binary outcome equals 1 if the deterministic validator would accept the final candidate for deployment and 0 otherwise. The preregistered missingness\_plan excluded a paired unit only when both paired runs were technical failures (model API failure, verifier timeout, or parser error). Execution accounting: planned\_episode\_count = 300, completed\_episode\_count = 282, failed\_episode\_count = 18, complete\_pairs = 141 (analysis/missingness\_summary.csv). Both-missing pairs = 9; provider-call audit: hosted\_model\_call\_event\_count = 156, completed = 147, failed = 9 (analysis/deterministic\_reconciliation.json). All exclusions and failure reasons are archived in execution/execution\_log.jsonl and analysis/missingness\_summary.csv.

Statistical-analysis procedures (artifact-backed). The preregistered primary test is the exact two-sided McNemar test on paired binary outcomes; contingency counts used by the test are archived at analysis/paired\_contingency\_table.csv (n11 = 135, n10 = 6, n01 = 0, n00 = 0). Paired effect estimation is presented as the absolute difference in success rates (guarded - baseline) with a bootstrap-percentile 95\% CI obtained from 10,000 resamples (bootstrap\_seed = 1729); bootstrap parameters and logs are archived in analysis/analysis\_log.jsonl. Secondary analyses (exploratory) were planned for subgroup effects by mutation class (string-level, token-injection, slot-override) and for a benign-control false-positive-blocking rate; the benign-control (M = 50) was not executed in this run and therefore no empirical false-positive-blocking rates are reported.

Determinism, provenance, and reproducibility notes. All run artifacts required to reproduce the analysis are archived and referenced above (execution/* and analysis/*). The single protocol deviation (temperature = null in execution/model\_configuration.json) is disclosed; because the pipeline used shared\_initial\_candidate semantics and provider-call accounting shows a single shared generation per pair (stage\_counts.shared\_initial\_generation = 150), paired comparability is preserved for this run. For deterministic reproduction we recommend explicitly setting temperature = 0 and re-running the archived execution/task\_manifest.jsonl with identical orchestration to produce bit-identical provider traces under the same provider semantics. Representative artifact paths for forensic verification: execution/task\_manifest.jsonl, execution/responses/task-000001-shared-initial.txt, execution/scoring/task-000001-baseline.json, execution/scoring/task-000001-guarded.json, analysis/deterministic\_reconciliation.json, analysis/paired\_contingency\_table.csv, analysis/analysis\_log.jsonl.

Pre-registered safeguards and contamination controls. The preregistration included contamination controls (artifact hashing and embargo rules for hazardous artifacts) and explicit exclusion rules for sanitized benign intents with excessive slot-drop (>0.3). Contamination summaries and integrity checks are archived in analysis/contamination\_summary.csv and the execution/execution\_manifest.json. Any artifact-flagging and masking procedures applied during archival are recorded; flagged hazardous raw artifacts are embargoed per the preregistered contamination policy.

Summary of methodological guarantees and limitations. Methodologically, the design enforces a paired comparison isolating downstream guarded-pipeline effects by reusing a single shared\_initial\_candidate per pair and limiting repairs to a single bounded attempt. Auditability is provided by per-episode validator traces, provider-call traces, and deterministic reconciliation artifacts. Limitations inherent in the design (single-model, synthetic tasks, validator coverage, and the protocol deviation noted above) are described in Threats to Validity and Limitations; these do not alter the integrity of the reported artifact-grounded paired comparison but bound external generalization.

\section{Results}
Execution and counts. The frozen run produced planned episodes = 300, completed episodes = 282, failed episodes = 18, and model provider calls used = 156 (execution/execution\_manifest.json). The primary confirmatory analysis used 141 complete pairs (analysis/missingness\_summary.csv). Condition summary (analysis/condition\_summary.csv) shows baseline\_successes = 135/141 (0.9574468085) and guarded\_successes = 141/141 (1.0).

Paired contingency and inference. The archived paired contingency (analysis/paired\_contingency\_table.csv) is n11 = 135, n10 = 6 (guarded=1, baseline=0), n01 = 0, n00 = 0. An exact two-sided McNemar test on the observed paired outcomes yields p = 0.03125. The point estimate (guarded - baseline) = 0.04255319148936165. The 95\% bootstrap percentile CI (10,000 resamples, seed = 1729) is [0.014184397163120567, 0.07801418439716312]; bootstrap parameters and log entries are archived in analysis/analysis\_log.jsonl.

Discordant-case diagnostics. The six discordant n10 pairs are fully traceable in the archive and are associated with repair-stage activity: analysis/deterministic\_reconciliation.json records repair = 6 and the corresponding per-episode scoring artifacts contain validator before/after snapshots showing sanitizer and/or permitted repair changes that converted a rejected initial draft into an accepted final candidate. Representative artifacts include execution/responses/task-000001-shared-initial.txt, execution/scoring/task-000001-baseline.json and execution/scoring/task-000001-guarded.json; these artifacts explicitly show the identical shared\_initial\_candidate referenced by both scoring artifacts and the validator's validation\_before and validation\_after objects.

Missingness sensitivity and unexecuted items. Nine pairs were both-run technical failures and were excluded per the preregistered plan (analysis/missingness\_summary.csv: both\_missing\_pairs = 9). Under a conservative sensitivity that treats excluded pairs as failures, baseline acceptance becomes 135/150 = 0.9000 and guarded acceptance 141/150 = 0.9400. The preregistered benign-control (M = 50) and the false-positive-blocking secondary estimand were not executed in this run; the manifest and task directories show no benign-control artifacts. Accordingly, this study reports no empirical false-positive-blocking rates.

\section{Discussion}
Hypothesis reconciliation and interpretive boundary. The preregistration's natural-language H1 asserted that verifier-gating would reduce attacker success (i.e., fewer unsafe configurations accepted). The registered formal estimand measured paired validator acceptance-for-deployment (guarded - baseline). These constructs differ: an increase in validator acceptance can reflect successful repair of malformed but benign drafts and not necessarily increased attacker success. Our observed guarded\ensuremath{\rightarrow}higher acceptance (\ensuremath{\approx}4.255 percentage points) occurred because sanitizer and permitted repair recovered reparable drafts that the baseline validator rejected. Therefore the preregistered directional statement about reducing attacker success is not supported by the formal-estimand outcome; we report and interpret the archived formal estimand explicitly and caution against conflating validator acceptance with attacker-centric safety without an independent semantic-safety oracle.

Operational implications and repair--gate co-design. Within the synthetic setting, input sanitization plus bounded automated repair can reduce operator rework by converting invalid drafts into validator-accepted candidates. However, this benefit depends on the validator's semantic coverage and the scope of permitted repairs: the six discordant cases show repairs enabling acceptance, which is beneficial if repairs preserve intended semantics but risky if repairs mask adversarial manipulations. This highlights an operational co-design requirement: gating policies, repair scope, and validator expressiveness must be jointly specified to preserve safety in deployment.

Practical trade-offs illustrated by the archive. The artifact evidence (analysis/deterministic\_reconciliation.json, analysis/paired\_contingency\_table.csv, and per-episode scoring artifacts such as execution/scoring/task-000001-guarded.json) shows precisely how small numbers of permitted repairs materially changed final acceptance. That pattern exposes a predictable trade-off: stricter gating (reject-on-violation) reduces the chance of a repaired but semantically altered candidate being accepted, at the cost of more operator interventions; more permissive policies with bounded repairs reduce operator workload but raise the need for stronger semantic checks or human review on repaired cases. These trade-offs are operational, measurable with the archived logs (repair stage\_count = 6, hosted\_model\_call\_event\_count = 156) and must be chosen to reflect acceptable risk budgets in deployment.

Failure modes of generate--validate--repair loops. The archived validator traces document three related failure modes relevant to operational design: (1) reparable syntactic/ordering errors where repair restores intent-preserving structure (the desirable class observed for the six discordants); (2) repair that fills missing fields or normalizes values in ways that change semantics subtly (requires semantic-safety or reachability checks beyond the current validator); and (3) cases of technical failure or timeouts (18/300 episodes, both-run failures = 9) where neither repair nor gating completed, reducing effective sample size (analysis/missingness\_summary.csv). Operational pipelines should instrument and surface these distinct failure categories to operators and governance audits using the same per-episode trace primitives present in our archive.

Reproducibility and the protocol deviation. The unset runtime sampling temperature (temperature = null in execution/model\_configuration.json) is a recorded protocol deviation. Its practical implication depends on provider default behavior; because the run used shared\_initial\_candidate semantics (stage\_counts.shared\_initial\_generation = 150) and provider-call accounting confirms the observed call pattern (analysis/deterministic\_reconciliation.json), paired comparability for the formal estimand is preserved in this execution. Nevertheless, for deterministic reproduction we recommend setting temperature = 0 explicitly (execution/model\_configuration.json) and re-running the archived task manifest; the archive provides all artifacts necessary to compare outcomes deterministically (execution/responses/*, execution/scoring/*, analysis/*).

Governance, auditing, and measurement recommendations grounded in archived evidence. The run demonstrates the value of (a) detailed per-episode validator traces (validation\_before/after objects in scoring artifacts) to audit repairs; (b) explicit repair audit logs and provenance (present for repaired episodes) so that operator review can target only changed cases; and (c) a preregistered benign-control to quantify false-positive-blocking (not executed here). The archive's contamination and missingness summaries (analysis/contamination\_summary.csv is empty; analysis/missingness\_summary.csv reports counts) show these governance metrics can be computed automatically from the same instrumentation used in this study.

Limits to operational generalization. The discussion above is tethered to the archive: a single LLM (declared gpt-5-mini), a single deterministic validator implementation, and synthetic templates. These constraints mean the measured acceptance-rate change is a within-bundle operational observation, not a claim of cross-vendor or production robustness. The ceiling effect (high baseline acceptance) reduces practical headroom for improvement; this is visible in analysis/condition\_summary.csv and motivates future runs that (1) execute the preregistered benign-control to measure false positives, (2) adopt attacker-centric semantic-safety or reachability oracles in addition to syntactic validators, and (3) diversify model and validator families before making operational deployment claims.

Concluding interpretive note. The archived data concretely show that bounded verifier-gated repair can increase validator acceptance by recovering reparable drafts; whether that outcome improves or harms security depends on whether repairs preserve intended semantics---a property not fully covered by the present deterministic validator and therefore not directly measured here. The empirical artifacts (per-episode scoring traces, repair counts, and paired contingency outputs) provide the necessary operational observables to design follow-up confirmations that target attacker-centric estimands and benign-control measurements.

\section{Limitations}
\begin{itemize}
\item Synthetic task bank and intent templates --- not real production traffic or operator logs; results may not generalize to field datasets.
\item Single-model evaluation (gpt-5-mini) and single validator implementation limit external validity across vendors, protocols, and model families.
\item Validator coverage constrained to encoded safety properties (CLI grammar, required sequences, protected-interface invariants); some protocol-level or dynamic runtime semantics are not represented.
\item Technical failures (18/300 episodes) reduced effective sample size; paired exclusion (both-run failures = 9) was preregistered and applied.
\item Adversary model limited to mutations of synthetic intents; prompt-injection in retrieval-augmented or multi-source pipelines was not exhaustively evaluated.
\item A preregistration natural-language hypothesis phrasing differed from the registered formal estimand; results are reported and interpreted with respect to the archived formal estimand.
\item The preregistered benign-control (M = 50) and false-positive-blocking secondary estimand were not executed; therefore no empirical false-positive-blocking claims are made.
\end{itemize}

\section{Conclusion}
A preregistered, artifact-backed paired experiment on a synthetic adversarial intent bank found that template-based input sanitization combined with deterministic verifier-gating and at-most-one automated repair produced a small but statistically detectable absolute increase in deterministic-validator-accepted LLM-generated configurations under the archived formal estimand (paired difference \ensuremath{\approx} 0.04255; p = 0.03125; 95\% CI [0.01418, 0.07801]). Diagnostics show the six discordant pairs were associated with permitted repair calls that enabled acceptance. Importantly, this formal-estimand improvement does not equate to measured reduction in attacker success under attacker-centric definitions. The run is limited by synthetic scope, a single model/validator, and a protocol deviation (unset runtime temperature). Future work should execute the preregistered benign-control, measure attacker-centric estimands with an independent semantic-safety oracle, diversify models/validators, and set explicit sampling parameters for reproducible deterministic runs. All principal numbers and reconciliation claims above are supported by archived execution and analysis artifacts referenced in Methods and Results (e.g., execution/task\_manifest.jsonl; execution/responses/task-000001-shared-initial.txt; execution/scoring/task-000001-baseline.json; execution/scoring/task-000001-guarded.json; analysis/paired\_contingency\_table.csv; analysis/condition\_summary.csv).

\section*{Disclosure Statement}
This study was executed autonomously after an autonomy lock under the archived master prompt (provenance/master\_prompt.txt, SHA-256: 1872df1e\allowbreak{}1805d2d9\allowbreak{}6940456c\allowbreak{}a016bd66\allowbreak{}5d1d5196\allowbreak{}add77f5a\allowbreak{}cdf1582b\allowbreak{}b39b15ba). Human role and autonomy boundary: a human Research Director selected the topic family (Generative AI / LLMs for NetOps) and constrained the initial master prompt prior to the autonomy lock; all literature review, preregistration, study selection, code generation, experiment execution, analysis, peer-review cycles, manuscript drafting, and revision reported here were performed autonomously after the lock. Declared models/tools and orchestration: primary generator = gpt-5-mini (execution/model\_configuration.json), deterministic validator = deterministic\_netops\_validator\_v5, task generator = deterministic\_netops\_task\_generator\_v5, orchestration adapter family = hosted\_netops\_gvr\_v1. Preregistration: study cand-2026-03 was preregistered and archived prior to observing outcomes (preregistration\_sha256: ba6ed25e\allowbreak{}84c7240f\allowbreak{}be94f2d5\allowbreak{}0b3e778f\allowbreak{}bfffd810\allowbreak{}ff2db632\allowbreak{}df8946e8\allowbreak{}5f0f1373). Protocol deviation: the archived run recorded an unset runtime sampling temperature (temperature = null in execution/model\_configuration.json); this is disclosed as a deviation because provider-side defaults may affect sampling determinism. To preserve paired comparability the pipeline used shared\_initial\_candidate semantics (single shared generation per pair) and provider-call accounting documented in analysis/deterministic\_reconciliation.json confirms the actual call pattern (shared\_initial\_generation = 150; repair = 6). Key archived artifact references (examples): execution/task\_manifest.jsonl; execution/responses/task-000001-shared-initial.txt; execution/scoring/task-000001-baseline.json; execution/scoring/task-000001-guarded.json; analysis/deterministic\_reconciliation.json; analysis/paired\_contingency\_table.csv; analysis/condition\_summary.csv; analysis/analysis\_log.jsonl; execution/model\_configuration.json; execution/execution\_manifest.json. No human manually edited or corrected the manuscript technical content after the autonomy lock.

\begin{thebibliography}{99}
\bibitem{ref1} http://citeseerx.ist.psu.edu/viewdoc/summary?doi=10.1.1.300.8659
\bibitem{ref2} https://doi.org/10.1145/2342441.2342452
\bibitem{ref3} https://doi.org/10.1145/3626111.3628194
\bibitem{ref4} https://doi.org/10.1145/3605764.3623985
\bibitem{ref5} http://arxiv.org/abs/2302.04761
\bibitem{ref6} https://doi.org/10.1145/3656296
\bibitem{ref7} Hao Zhou, Chengming Hu, Ye Yuan, Yufei Cui, Yili Jin, Can Chen, Haolun Wu, Dun Yuan, Li Jiang, Di Wu, Xue Liu, Jianzhong Zhang, Xianbin Wang, Jiangchuan Liu, "Large Language Model (LLM) for Telecommunications: A Comprehensive Survey on Principles, Key Techniques, and Opportunities", 2024, 10.1109/comst.2024.3465447
\bibitem{ref8} Lingzhe Zhang, Tong Jia, Mengxi Jia, Yifan Wu, Aiwei Liu, Yong Yang, Zhonghai Wu, Xuming Hu, Philip S. Yu, Ying Li, "A Survey of AIOps in the Era of Large Language Models", 2025, 10.1145/3746635
\bibitem{ref9} Rajdeep Mondal, Alan Tang, Ryan Beckett, Todd Millstein, George Varghese, "What do LLMs need to Synthesize Correct Router Configurations?", 2023, 10.1145/3626111.3628194
\bibitem{ref10} Kai Greshake, Sahar Abdelnabi, Shailesh Mishra, Christoph Endres, Thorsten Holz, Mario Fritz, "Not What You've Signed Up For: Compromising Real-World LLM-Integrated Applications with Indirect Prompt Injection", 2023, 10.1145/3605764.3623985
\bibitem{ref11} Yupeng Chang, Xu Wang, Jindong Wang, Yuan-Hsuan Wu, Linyi Yang, Kaijie Zhu, Hao Chen, Xiaoyuan Yi, Cunxiang Wang, Yidong Wang, Wei Ye, Yue Zhang, Yi Chang, Philip S. Yu, Qiang Yang, Xing Xie, "A Survey on Evaluation of Large Language Models", 2023, 10.48550/arxiv.2307.03109
\bibitem{ref12} Timo Schick, Jane Dwivedi-Yu, Roberto Dessì, Roberta Răileanu, María Lomelí, Luke Zettlemoyer, Nicola Cancedda, Thomas Scialom, "Toolformer: Language Models Can Teach Themselves to Use Tools", 2023, 10.48550/arxiv.2302.04761
\bibitem{ref13} Changjie Wang, Mariano Scazzariello, Alireza Farshin, Simone Ferlin, Dejan Kostić, Marco Chiesa, "NetConfEval: Can LLMs Facilitate Network Configuration?", 2024, 10.1145/3656296
\bibitem{ref14} Peyman Kazemian, George Varghese, Nick McKeown, "Header space analysis: static checking for networks", 2012
\end{thebibliography}

\end{document}
